\clearpage

\section{Übung 2: Abfragen über mehrere Tabellen}

\subsection{Aufgabe 1: Tabellen verknüpfen \tagPflicht}
\noindent \emph{Neue Sprachmittel: Verknüpfung über Primär- und Fremdschlüssel (\texttt{INNER JOIN ... ON}), Tabellen-Aliase (\texttt{AS}).}

Führe auf der Datenbank \myUserInput{haustier} die folgenden Abfragen durch:

\begin{compactenum}[a)]
\item Zeige ohne Verwendung von \lstinline|JOIN| oder \lstinline|WHERE| eine
Tabelle mit allen Besitzern (Nachname und Vorname) und Tieren (Name und
Tierart) an. Was für eine Tabelle wird erstellt?
\item Zeige eine Liste mit allen Besitzern und den zu ihnen gehörenden Tieren
an. Es sollen alle Spalten aus \myUserInput{besitzer} und \myUserInput{tier}
angezeigt werden.
\item Zeige eine Tabelle mit allen Besitzern (Vor- und Nachname) und ihren
Tieren (Name und Tierart) an. Sortiere die Liste absteigend nach dem Namen der
Tiere.
\item Zeige eine Tabelle mit allen Besitzern (Nachname und Vorname) und ihren
Tieren (Name und Tierart) an. Sortiere die Liste in aufsteigender
alphabetischer Reihenfolge zunächst nach dem Nachnamen und dann nach dem
Vornamen.
\item Zeige eine Tabelle mit allen Besitzern (Vor- und Nachname) und ihren
lebendigen Tieren an (Name und Tierart). Sortiere die Liste nach der Tierart.
\item Zeige alle Tiere von Mirco Sandelmann an (Name, Tierart und
\myUserInput{lebendig}).
\item Wähle alle Besitzer von Hunden aus. Zeige Vor- und Nachname der Besitzer
sowie Name, Geburtstag und Todestag des Hundes an.
\item Liste alle Besitzer von Maja auf (Vor- und Nachname).
\item Zähle die Anzahl der Besitzer von Blacky.
\item Erstelle eine Liste, die angibt, wie viele Tiere jede einzelne Person
besitzt (Ausgabe: Anzahl, Vor- und Nachname). Sortiere die Liste in
umgekehrter alphabetischer Reihenfolge nach dem Nachnamen.
\item Zeige alle Besitzer mit Vor- und Nachnamen an, die zwei oder mehr Tiere
besitzen.
\item Zähle die Anzahl der Besitzer, die einen Hund oder eine Katze besitzen
(achte darauf, Besitzer mit mehreren Tieren nicht doppelt zu zählen!).
\item Liste die vollständigen Daten aller Besitzer auf, deren Telefonnummer
nicht bekannt ist.
\item Liste Straße, PLZ und Ort aller Besitzer auf. Dabei soll jede Adresse nur
einmal ausgegeben werden.
\end{compactenum}

\subsection{Aufgabe 2: \texttt{LEFT JOIN} und Unterabfragen \tagPflicht}
\noindent \emph{Neue Sprachmittel: Äußere Verknüpfungen (\texttt{LEFT JOIN ... ON} mit \texttt{IS NULL}) sowie geschachtelte Unterabfragen (\texttt{Subqueries} mit \texttt{=} und \texttt{IN}).}

Führe auf der Datenbank \myUserInput{haustier} die folgenden weiterführenden
Abfragen durch:

\begin{compactenum}[a)]
\item Liste die Personen, die momentan kein Tier besitzen, mit ihrem Vor- und
Nachnamen auf.
\item Liste die Tiere mit Namen auf, die keinen Besitzer haben.
\item Liste alle Personen, die mit Anka Anderson zusammenwohnen (d.\,h.\ alle,
die dieselbe Straße haben), mit ihren vollständigen Daten auf. Anka selbst
darf auch in der Liste erscheinen.
\item Liste alle Tiere mit Namen auf, die am selben Tag geboren wurden wie Maja.
Maja selbst soll nicht in der Liste erscheinen.
\item Liste die Namen der Tiere auf, die am selben Tag geboren wurden wie Maja,
sowie den Vor- und Nachnamen ihrer Besitzer. Maja selbst soll nicht in der
Liste erscheinen.
\item Liste alle Tiere mit Namen und Tierart auf, die denselben Besitzer haben
wie Hasso. Hasso darf ebenfalls in der Liste erscheinen, aber kein Tier soll
doppelt genannt werden.
\end{compactenum}